\section{The quotient and its lifted automorphisms}
\label{sec:quotient-lift}

Let
\[
\pi\colon G60\longrightarrow G30=G60/\langle a\rangle
\]
be the quotient by the canonical deck involution. This section determines
the automorphism group of the quotient and constructs its full lift to
the native graph.

\subsection{The quotient graph}

The quotient $G30$ has thirty vertices, one for each $a$-fiber of $G60$.
The exact quotient certificate records the fiber partition, the quotient
edge set, and the projection map $\pi$.

The quotient automorphism group has order
\[
\left|\operatorname{Aut}(G30)\right|=240
\]
and abstract structure
\[
\operatorname{Aut}(G30)\cong S_5\times C_2.
\]

The first factor reflects the canonical $S_5$ symmetry visible through
the Petersen and line-graph architecture. The second factor is an
additional central involution in the quotient action.

\subsection{The lifting problem}

For
\[
g\in\operatorname{Aut}(G30),
\]
a lift of $g$ is an automorphism
\[
\widetilde g\in\operatorname{Aut}(G60)
\]
satisfying
\[
\pi\circ\widetilde g=g\circ\pi.
\]

Because $\pi$ is a two-fold cover with deck group
\[
\langle a\rangle\cong C_2,
\]
any two lifts of the same quotient automorphism differ by composition
with $a$. Thus, whenever a lift exists, the complete lift fiber is
\[
\{\widetilde g,\ a\widetilde g\}.
\]

Existence is not automatic for an arbitrary automorphism of an arbitrary
cover. It must be checked against the native edge law.

\subsection{Exact lifting theorem}

The quotient certificate and native edge set permit a direct lift test
for every element of $\operatorname{Aut}(G30)$. All $240$ quotient
automorphisms pass.

\begin{theorem}
\label{thm:two-lifts}
Every automorphism of $G30$ has exactly two lifts to $G60$.
Consequently, the complete lifted set has cardinality
\[
2\left|\operatorname{Aut}(G30)\right|
=
2\cdot240
=
480.
\]
\end{theorem}

\begin{proof}
Let $g\in\operatorname{Aut}(G30)$. The exact lift construction assigns
to each native vertex an image lying over the quotient vertex
$g(\pi(v))$. The resulting permutation is tested against the complete
native edge set and is an automorphism of $G60$.

Composition with the nontrivial deck transformation $a$ gives a second
lift. These lifts are distinct because $a$ is fixed-point-free.

Conversely, two lifts of the same quotient automorphism differ by a deck
transformation. Since the deck group is $\langle a\rangle$ of order $2$,
there are at most two lifts. Hence there are exactly two.
\end{proof}

Let
\[
L
=
\left\{
\widetilde g:
g\in\operatorname{Aut}(G30)
\right\}
\cup
\left\{
a\widetilde g:
g\in\operatorname{Aut}(G30)
\right\}.
\]
The lift construction is closed under permutation composition and
inversion, so
\[
L\leq\operatorname{Aut}(G60)
\]
with
\[
|L|=480.
\]

\subsection{The extension over the quotient group}

The lifted group fits into an exact sequence
\[
1
\longrightarrow
\langle a\rangle
\longrightarrow
L
\longrightarrow
\operatorname{Aut}(G30)
\longrightarrow
1.
\]

The kernel is central in $L$. Indeed, every lifted automorphism commutes
with the deck involution:
\[
\widetilde g\,a=a\,\widetilde g.
\]

The extension is not described merely by adjoining an independent
central $C_2$ to the quotient group. Its internal structure contains
both split and nonsplit behavior over natural index-two subgroups of
$\operatorname{Aut}(G30)$. This distinction will be used in
Section~\ref{sec:group-structure} to identify the full group.

At this stage we have proved only that
\[
L\leq\operatorname{Aut}(G60)
\qquad\text{and}\qquad
|L|=480.
\]
It remains possible, in principle, that the native graph has additional
automorphisms not preserving the $a$-fiber system. The next section
eliminates that possibility by an independent enumeration performed
directly on $G60$.
